\begin{table}[pos=htbp]
\caption{Security-feedback outcomes stratified by initial visible state. Final pass rates are Secure-Pass@1 (\%). F-only and F+S denote the two feedback conditions. Each branch makes exactly two revisions. Tasks can occur in multiple strata. F.Reg. counts hidden functional regression relative to the shared initial candidate. These are post-hoc descriptive strata.}
\centering\small
\begin{tabular}{lrrr rrr}
\toprule
Initial visible & Initials & Tasks & \shortstack{Secure-Pass@1\\F-only (\%)} & \shortstack{Secure-Pass@1\\F+S (\%)} & Wins/losses/ties & F.Reg./F+S \\
\midrule
F+S+ & 100 & 12 & 85.0 & 84.0 & 0/1/99 & 1 \\
F+S- & 18 & 4 & 0.0 & 16.7 & 3/0/15 & 2 \\
F-S+ & 50 & 7 & 28.0 & 26.0 & 1/2/47 & 0 \\
F-S- & 32 & 6 & 12.5 & 21.9 & 3/0/29 & 0 \\
\bottomrule
\end{tabular}
\end{table}
\begin{table}[pos=htbp]
\caption{Project-level B0-to-M outcomes by initial visible state. Repair and regression counts use the initial hidden state; a visible functional failure need not imply a hidden functional failure. Calls exclude the shared initial generation.}
\centering\small
\begin{tabular}{lrrrrrr}
\toprule
Initial visible & Initials & Tasks & J repairs & F.Reg. & \shortstack{Secure-Pass@1\\M (\%)} & New calls \\
\midrule
F+S+ & 27 & 6 & 0 & 0 & 85.2 & 0 \\
F+S- & 28 & 7 & 6 & 13 & 21.4 & 47 \\
F-S+ & 39 & 6 & 12 & 0 & 53.8 & 64 \\
F-S- & 6 & 1 & 0 & 3 & 0.0 & 12 \\
\bottomrule
\end{tabular}
\end{table}
\begin{table}[pos=htbp]
\caption{Project-level functional regression mechanisms. Categories describe the observed execution failures in all 16 B0-to-M regressions.}
\centering\small
\begin{tabular}{lrl}
\toprule
Mechanism & Count & Cases \\
\midrule
Response extraction / loading failure & 6 & EX02, EX08, EX10--11, EX15--16 \\
Unpickler used as context manager & 4 & EX04--07 \\
Template syntax or interpolation damage & 3 & EX12--14 \\
AST interpreter semantic errors & 2 & EX01, EX03 \\
Overrestrictive preference schema & 1 & EX09 \\
\bottomrule
\end{tabular}
\end{table}
\FloatBarrier
\paragraph{Case selection and evidence}
The artifact ledger includes all ten security-feedback joint disagreements, plus two additional functional regressions that tie on joint failure, and all 16 project-level B0-to-M functional regressions. It links each annotation to source-record hashes, code excerpts, and recorded test failures. Case classifications use machine-assisted source/log inspection. Automated checks verify completeness of case selection and correspondence between annotations and archived evidence. The analysis uses the archived generations and original outcome scores.

\paragraph{Security-feedback cases}
SF01--04 add a null-result guard after \texttt{safe\_join} to handle an exception on rejected traversal input. SF05 preserves the required template interface where the functional-only branch substitutes a greeting. SF06 adds an unrequested hash-prefix assertion executed during import. SF07--08 use \texttt{escape} without an import. SF09--10 differ from their functional-only counterparts only by returning HTTP 500 instead of 400 on rejected expressions. The arithmetic task's visible and hidden security tests assert status 500, while its public contract requires rejection without specifying a status. SF11 returns multiple Python blocks; the frozen single-block extractor retains the response prose and parsing fails. SF12 introduces unavailable \texttt{sympy}. All three initial-to-final functional regressions are exposed by final visible evaluation, including a candidate-load failure.

\paragraph{Project-level regressions}
EX01 resolves \texttt{input\_data} as a dictionary key instead of the variable binding; EX03 rejects string subscripts. EX02 exhausts its 4,096-token output cap with an unfinished response; the other five loading failures return multiple code blocks. EX04--07 raise a context-manager \texttt{TypeError} before unpickling benign or malicious data. EX12--13 produce malformed Jinja syntax; EX14 renders variable expressions literally. Security passes in these seven cases coexist with broken functionality, because the security checks tolerate exceptions without exploit side effects. EX09 restricts preferences to \texttt{theme}/\texttt{language}, rejecting the legitimate hidden \texttt{notifications} key; its visible functional example contains only the two allowed keys. Validation after \texttt{pickle.loads} also fails to prevent recorded malicious side effects.

Fifteen of the 16 regressions fail the final visible functional test and remain unresolved at the revision limit. EX09 passes the visible test but fails on a legitimate hidden input. The cases therefore distinguish incomplete repair of exposed failures from a gap in visible test coverage.

\paragraph{Opportunity, costs, and sensitivity}
Initial-state strata retain all 200 paired initials and all 100 project-level B0 initials. No stratum uses revised code or hidden outcomes for assignment. The accompanying CSV gives Pass@1, Secure@1, Secure-Pass@1, repairs, regressions, calls, and mean input/output tokens for each stratum; JSON additionally retains model breakdowns for the security-feedback experiment. The whole-task sensitivity analysis excludes all ten arithmetic-task pairs: functional-only and combined feedback score 54.21\% and 55.26\% joint success on the remaining 190 pairs. This post-hoc sensitivity analysis accompanies the full 20-task results.
